\section{同一个问题怎样跨核区比较}\label{sec:4-1}

在第~1~章中，原子核低能激发谱已被置于六个典型质量区的宏观坐标之中；第~2~章与第~3~章则进一步从微观壳模型、平均场自洽形变、对称性恢复投影以及转动、手征与摆动等集体动力学角度，奠定了多体波函数与唯象图像的数理基础。当分析视角从单个核素转向全局核素图谱时，跨核区比较往往成为检验模型泛化能力、揭示演化规律的最有力工具。然而，跨核区比较也是极易滋生虚假类比与概念滑移的区域。

最常见的认识论误区在于：仅凭某一核区发现了特定能带结构（如两条近简并带或交错低自旋带），便将质量数 $A$ 或单一唯象参数（如能级比 $E(2^+_2)/E(2^+_1)$、三轴形变角 $\gamma$ 的单一平衡值）作为排他性判据，孤立地宣布在另一区域也必定存在完全相同的物理模式。事实上，核多体系统的集体响应从来不是孤立几何形状的直接映射，而是价轨道占据、有效相互作用单极漂移、配对超流关联与自旋窗口内探针选择定则综合作用的宏观显现。本节旨在建立一套跨核区比较的物理判据与分析方法，明确哪些量子几何与多极代数是普适可迁移的，哪些微观输入与有效参数需要在新体系中重新检验。

\subsection{比较对象：内禀组态、实验室对称性与态的映射}\label{sec:4-1-1}

跨核区比较的首要前提是厘清``究竟在比较什么对象''。原子核作为自约束的有限量子多体系统，在没有外场作用的实验室参考系中，由强相互作用与电磁相互作用主导的哈密顿量在极高精度下严格具备空间旋转群 $\mathrm{SO}(3)$ 对称性、空间反演宇称 $\mathcal{P}$ 对称性以及时间反演 $\mathcal{T}$ 对称性（在此忽略微弱的弱相互作用宇称不守恒效应）。因此，实验室参考系中物理定态的基矢必定可以按确定的总角动量量子数 $I$、宇称 $\pi$ 及其磁量子数 $M$ 展开。在实际谱学实验中，测得的通常是具有确定能级位置的纯态或处于特定取向系综与统计混合的退极化布居。

自洽平均场模型中所讨论的轴对称长椭球、扁椭球或三轴椭球形状，本质上是内禀对称性破缺参考态 $|\Phi(\beta,\gamma)\rangle$ 的几何表象。内禀系中的非零四极形变参量 $(\beta, \gamma)$ 描述的是核内多核子空间密度分布的二阶各向异性关联，如第~2.4~节式~\eqref{eq:2-4-2}所引入的玻尔坐标。

然而，单一内禀态 $|\Phi(\beta,\gamma)\rangle$ 本身并非实验室物理本征态。利用角动量投影算符 $\hat{P}^I_{MK}$ 对单一内禀构型进行全欧拉角积分，所构造的仅是具备良好角动量量子数的投影试探态（trial states）：
\begin{equation}\label{eq:4-1-projection}
|IM\alpha\rangle = \sum_K g_K^\alpha \hat{P}^I_{MK} |\Phi\rangle = \sum_K g_K^\alpha \frac{2I+1}{8\pi^2} \int \mathrm{d}\Omega \, \mathcal{D}^{I*}_{MK}(\Omega) \hat{R}(\Omega) |\Phi\rangle.
\end{equation}
在物理真实的有限核系统中，实际定态通常是多个具有不同形变构型、配对自由度或准粒子激发的内禀态的多体量子叠加（例如生成坐标法 GCM 或多参考能量泛函理论所描述的构型混合网络）。

这意味着，跨核区比较的对象绝不能停留在某个未投影平均场势能面（PES）局域极小点的坐标数值上，而必须落实到经对称性恢复后具有确定自旋与宇称的物理态 $|I^\pi\rangle$、这些态之间的跃迁网络以及波函数在形变构型空间上的概率分布。在波函数分布极其弥散的大振幅软核中，实验室一体电荷密度对于 $0^+$ 基态处处呈现球对称，唯有通过高阶两体密度关联与电磁多极约化矩阵元，方能提取内部形状的客观物理实在性\cite{poves-2020-limits-assigning-nuclear-shape}。

\subsection{微观驱动与动力学响应：轨道占据、单极演化与配对关联}\label{sec:4-1-2}

比较本文所选的核与自旋区间时，一个重要差别是费米面附近的轨道及其占据。以下是相关案例中的轨道线索；同一质量数附近仍可有不同的质子数、中子数与组态。
\begin{itemize}
\item 在 $A\approx 60$ 与 $80$ 区域，质子与中子费米面主要受 $1f_{5/2}, 2p_{3/2}, 2p_{1/2}$ 以及下落的 $1g_{9/2}$ 独特宇称轨道主导；
\item 在 $A\approx 130$ 区域，中子接近壳层中上部，质子位于准空穴或下部，负宇称独特轨道 $1h_{11/2}$ 的驱动效应成为主导；
\item 在 $A\approx 160$ 稀土重区，强形变长椭球势阱将 $1i_{13/2}$ 轨道大幅劈裂；
\item 在 $A\approx 190$ 区域，费米面则深埋于 $1i_{13/2}$ 中子空穴与 $1h_{9/2}, 2f_{7/2}$ 混合轨道之中。
\end{itemize}

在给定高 $j$ 壳层与近轴对称参考中，不同 $|\Omega|$ 分量具有不同的四极矩，粒子与相对于满壳定义的空穴也可给出相反贡献。费米面移动改变这些分量的占据，并通过残余相互作用与核芯极化反馈改变形变倾向。但单个 $K$ 或 $\Omega$ 标签不能独自决定长椭、扁椭或三轴形状；还需核对具体 Nilsson 能级随形变的变化、配对和相互作用。第~2.2~节的占据反馈提供了这种比较的微观出发点\cite{book-bohr-mottelson-nuclear-structure-vol2}。

更深层次的微观驱动还来自于多核子相互作用中的单极效应（monopole interaction）。正如第~2.2~节式~\eqref{eq:2-2-1}所展开的单极相互作用分解与有效单粒子能量（ESPE）演化框架所示，随着价核子在不同组态间的填充或跃迁，核子间张量力与自旋-同位旋交换相互作用引起有效能隙的动态漂移\cite{doctoral-r1-g28-otsuka-2023-type-ii-shell-evolution}。在跨区比较中，这种单极演化逻辑展现出鲜明的核区特性：例如在 $A\approx 100$ 区域，正是质子从 $2p, 1f$ 轨道激发至 $1g_{9/2}$ 时，通过强烈的单极吸引作用诱发了中子 $1g_{7/2}$ 与 $1h_{11/2}$ 轨道的突变性下移，形成了所谓的第二类壳演化（Type-II shell evolution）与形变自强化反馈\cite{doctoral-r1-g28-togashi-2016-zr-mcsm-qpt}；而在高质子数区，这种单极诱导相变的临界核子数则发生显著位移。

配对关联通过改变占据、准粒子能量和响应矩阵元调制转动惯量。在指定平均场的比较中，配对常使低频响应相对无配对参考减小；这不构成各轴惯量与经典刚体值之间的普遍不等式。第~2.6~节区分了冻结场与含残余响应的实现。不同核或组态的单粒子谱、配对和顺排过程不同，跨区比较时应重新检验这些输入。

\subsection{形变空间测度与高阶不变量：$\gamma$ 刚性、$\gamma$ 软性与分布宽度}\label{sec:4-1-3}

三轴自由度是跨区比较中最易引起混淆的概念之一。在唯象模型中，人们常将 Davydov-Filippov 刚性三轴转子模型与 Wilets-Jean $\gamma$ 不稳定模型进行对照：
\begin{itemize}
\item \textbf{刚性三轴假设}：设想势能面在非零三轴平衡角（如 $\gamma \approx 20^\circ\text{--}30^\circ$）处具有较深而局域的极小，波函数集中于极小点近旁；
\item \textbf{$\gamma$ 软性/不稳定假设}：势能面沿 $\gamma$ 方向展现宽平坦特性，核波函数在整个三轴区间广泛弥散，不存在静态固定的三轴取向。
\end{itemize}

指定的理想刚性转子与 $\gamma$ 不稳定模型可以给出不同的能级和跃迁预言，应逐项比较。在更一般的势能与惯量可调模型中，少数能级对反问题的约束仍可能不足。因此不能把某个能比当作形状分布的唯一测量，也不能把拟合的非唯一性误写成两个理想极限的完整能谱相同。

第~2.1、2.4~节给出了 Kumar--Cline 不变量的实验室张量定义\cite{cline-1986-nuclear-shapes-coulomb-excitation}。标量求和不需预先选定转子势，几何解释则仍遵循所声明的坐标、算符与截断条件。对 $0^+$ 态，沿用同一归一可写：
\begin{equation}\label{eq:4-1-q2-invariant}
P_2\equiv q_2=\langle Q^2\rangle=\sqrt5\langle0_i^+|[\mathcal M(E2)\times\mathcal M(E2)]_0|0_i^+\rangle=\sum_t|\langle2_t^+\Vert\mathcal M(E2)\Vert0_i^+\rangle|^2,
\end{equation}
\begin{equation}\label{eq:4-1-q3-invariant}
P_3\equiv q_3=\langle Q^3\cos3\gamma\rangle=-\sqrt{\frac{35}{2}}\langle0_i^+|[\mathcal M(E2)\times[\mathcal M(E2)\times\mathcal M(E2)]_2]_0|0_i^+\rangle.
\end{equation}
这里的几何解释沿用第~2.1、2.4~节的坐标与算符条件；原作形状角 $\delta$ 的记号须与本文辐射混合比分开。更高阶不变量可约束相应算符矩的涨落，但从这些矩到连续角度分布仍需额外假设，不能仅由四阶四极矩直接读出角度方差。

然而必须保持审慎的证据边界：低阶几何不变量 $\langle Q^2 \rangle$ 与 $\langle Q^3 \cos 3\delta \rangle$ 仅提供了宏观形状电荷分布的二阶与三阶矩重心，并不能唯一反演核多体波函数在 $(\beta,\gamma)$ 连续坐标空间上的全概率分布密度。更重要的是，实验测得的 $E2$ 矩阵元网络总是有限截断的，高能中间态的缺失与矩阵元相对符号的不确定性会引入系统误差。因此，分散度参数能够为检验多体模型波函数提供关键约束，但绝不能将其简化为某种机械数值门槛去绝对化地判定刚性或软核。

\subsection{实验探针与观测选择定则：自旋窗口、激发机制与电磁张量算符}\label{sec:4-1-4}

跨区比较结论的可信度，高度受制于所采用实验探针的布居物理机制及其自旋-宇称选择定则。不同实验手段在核素图上所能触及的激发能量与自旋窗口存在显著差异：
\begin{enumerate}
\item \textbf{低自旋、高精度绝对跃迁探针}：
   低能库仑激发、$\beta$ 衰变、中子俘获与NRF具有不同的形成机制和选择规则。中子俘获还涉及复合核，不能统称为纯电磁微扰探针。可达自旋由实际入射道、母态及级联决定；相应矩阵元或率的提取还需效率、分支、寿命或反应模型等输入，不能给这一组方法规定统一的低自旋上限。
\item \textbf{中高自旋在束反应探针}：
   包括重离子熔合蒸发反应（Heavy-Ion Fusion-Evaporation）、深部非弹性碰撞（Deep Inelastic Collisions）、多核子转移反应以及自发/诱发裂变碎块谱学。这类探针通过大角动量输入能够有效布居高自旋转动带，是探索高自旋带交叉、磁转动剪刀带、手征双重带及摆动激发的有力工具。然而，必须注意不同反应机制的角动量输入截面与退激统计分布截然不同；且多步库仑激发在重区强形变核中同样可以达到 $I \sim 20\text{--}30$ 的高自旋态。绝对跃迁强度的测定不仅依赖反冲距离法（RDDS）和多普勒频移衰减法（DSAM），还广泛依赖于电子学超快计时（如 $\mathrm{LaBr}_3(\mathrm{Ce})$ 阵列）、延迟符合、转化电子谱学以及 NRF 共振荧光线宽分析等互补技术。
\end{enumerate}

不同观测量约束波函数的不同侧面，其区分能力取决于待检验解释与实验误差，不存在普遍的信息量排名：
\begin{itemize}
\item \textbf{能级能量 $E(I)$}：约束哈密顿量与态混合，有时对交叉和能隙尤为敏感；当多个模型已用这些能量调参时，还需留出观测检验；
\item \textbf{跃迁强度 $B(E2)$ 与静四极矩 $Q_s$}：前者含非对角约化矩阵元的模平方和初态统计权，后者约束对角矩阵元；二者提供互补的电荷四极信息；
\item \textbf{$\rho^2(E0)$、同位素与异能态位移}：$E0$约束同核态间的径向跃迁矩阵元，位移比较相应电荷半径期望值。结合指定构型模型可限制混合与半径差，但不由单一量唯一重建形状；
\item \textbf{单核子转移反应与单粒子谱强}：实验直接测量的物理量是反应微分截面 $\mathrm{d}\sigma/\mathrm{d}\Omega$。提取单粒子光谱因子 $S_j$ 需要引入光学势畸变波波恩近似（DWBA）或绝热转移模型，其数值不仅依赖于反应动力学势的选择，还受限于实验激发能量分辨率以及截断窗口内的碎片求和；正如第~2.2~节所讨论的，单核子转移截面反映的是核多体波函数与单粒子试探算符的重叠，不能将其误解为裸算符的一体直接测量。
\end{itemize}

\subsection{可迁移规律与需重新检验条件的物理判据：六区比较总表}\label{sec:4-1-5}

基于上述微观机制与实验物理剖析，在跨质量区核结构比较中，必须严格区分普适可迁移的理论法则与不可无条件迁移的局域输入：
\begin{itemize}
\item \textbf{可迁移的物理法则}：
  \begin{enumerate}
  \item 空间旋转与反演对称性的量子几何代数：Wigner-Eckart 定理的约化张量定义、Clebsch-Gordan 系数相位与 Racah 重耦合代数\cite{book-suhonen-from-nucleons-to-nucleus}；
  \item 电磁辐射场多极展开与跃迁几率的基本微观定义；
  \item 模型无关的 Kumar-Cline 四极不变量理论框架与零动量标量积完备求和规则\cite{cline-1986-nuclear-shapes-coulomb-excitation}；
  \item 态混合与微观组态干涉的代数分析方法。
  \end{enumerate}
\item \textbf{不可直接迁移的局域物理条件}：
  \begin{enumerate}
  \item 具体单粒子轨道的有效单粒子能量（ESPE）与局部主壳能隙；
  \item 价空间截断所引入的模型空间有效电荷（$e_p^{\mathrm{eff}}, e_n^{\mathrm{eff}}$）与有效弱轴矢流重整化因子 $g_A^{\mathrm{eff}}$；
  \item 唯象模型（如 IBM、PRM、玻尔哈密顿量 5DCH）通过局部能谱拟合所得的势能形状参数、微观质量张量或刚体转动惯量；
  \item 某一特定核区因特定高 $j$ 独特轨道单粒子耦合产生的复杂带结构（如手征双重带），在另一核区轨道占据与三轴极小尚未独立确立前，绝不能仅凭能谱外貌相似直接类推存在相同模式。
  \end{enumerate}
\end{itemize}

根据此项比较判据，表~\ref{tab:4-1-six-regions-framework}系统梳理了六个典型质量区的核心物理驱动、主导费米面轨道、关键形变特征以及当前实验认知的确定性边界。在全书后续比较论述中，我们将严格遵守表中关于``明确测定''、``未见''、``未搜索''与``灵敏度不足''四类认识论边界的界定，绝不以跨区类比掩盖物理认知的真实空缺。

\begin{table}[htbp]
\centering
\small
\setlength{\tabcolsep}{2.5pt}
\caption{以本文所选案例为线索的六区比较；质量区不是互斥的模式分类，模型解释与直接报告量分别使用}\label{tab:4-1-six-regions-framework}
\begin{tabularx}{\textwidth}{p{1.5cm}p{3.5cm}p{3.6cm}X}
\toprule
质量区 & 主导轨道与占据特征 & 核心物理机制与现象 & 实验证据与认知边界 \\
\midrule
$A \approx 60$ & 所选Ge、Ni中的壳间激发与轨道占据 & 组态混合、次闭壳效应与不同形状的模型解释 & $^{64}\mathrm{Ge}$ 寿命与E2已测；$^{68}\mathrm{Ni}$ 的低能态寿命、E2与E0上限共同约束三形状解释 \\
\addlinespace
$A \approx 80$ & 所选Kr、Ge及奇奇核的价轨道占据与核芯响应 & 形状混合、四极涨落与取向候选须分核、分态比较 & Kr带符号矩约束混合；作者以Ge高阶不变量讨论刚性；As、Br问题见第4.5节 \\
\addlinespace
$A \approx 100$ & 所选Zr模型中的占据反馈及Mo、Ru的不同响应 & 急变、软性和多极响应可出现在不同态与能窗 & Zr寿命与模型比较；Mo低能态和相对分支；Ru报告偶极强度并结合模型提出八极解释 \\
\addlinespace
$A \approx 130$ & 所选质子轨道与中子粒子或空穴的组态组合 & 手征、摆动、进动及软核耦合的条件比较 & Pr、Ce等核已报告不同候选，须分别核对能带、绝对率和所用模型 \\
\addlinespace
$A \approx 160$ & 高 $j$ 轨道与多准粒子构型；Lu案例涉及 $\pi i_{13/2}$ & 大形变转动、高 $K$ 异能态及配对与顺排响应 & Lu的TSD带提供摆动比较基准；稀土高 $K$ 跃迁须按实际支路与禁戒度比较 \\
\addlinespace
$A \approx 190$ & 所选核接近闭壳区，粒子或空穴占据随核素而变 & 扁椭、三轴、形状混合及单粒子描述的交界 & Hg、Ir、Os、Au案例按自旋窗、组态、跃迁强度与竞争解释比较 \\
\bottomrule
\end{tabularx}
\end{table}
